\documentclass[sigconf,review,anonymous]{acmart}
%%
%% \BibTeX command to typeset BibTeX logo in the docs
\AtBeginDocument{%
  \providecommand\BibTeX{{%
    Bib\TeX}}}

%% Rights management information.  This information is sent to you
%% when you complete the rights form.  These commands have SAMPLE
%% values in them; it is your responsibility as an author to replace
%% the commands and values with those provided to you when you
%% complete the rights form.
\setcopyright{acmlicensed}
\copyrightyear{2018}
\acmYear{2018}
\acmDOI{XXXXXXX.XXXXXXX}
%% These commands are for a PROCEEDINGS abstract or paper.
\acmConference[Conference acronym 'XX]{Make sure to enter the correct
  conference title from your rights confirmation email}{June 03--05,
  2018}{Woodstock, NY}
%%
%%  Uncomment \acmBooktitle if the title of the proceedings is different
%%  from ``Proceedings of ...''!
%%
%%\acmBooktitle{Woodstock '18: ACM Symposium on Neural Gaze Detection,
%%  June 03--05, 2018, Woodstock, NY}
\acmISBN{978-1-4503-XXXX-X/2018/06}


\newcommand{\dk}[1]{\textcolor{teal}{(DK: #1)}}
\newcommand{\mf}[1]{\textcolor{red}{(MF: #1)}}
\newcommand{\mg}[1]{\textcolor{orange}{(MG: #1)}}
\newcommand{\yc}[1]{\textcolor{blue}{(YC: #1)}}


%%
%% Submission ID.
%% Use this when submitting an article to a sponsored event. You'll
%% receive a unique submission ID from the organizers
%% of the event, and this ID should be used as the parameter to this command.
%%\acmSubmissionID{123-A56-BU3}

%%
%% For managing citations, it is recommended to use bibliography
%% files in BibTeX format.
%%
%% You can then either use BibTeX with the ACM-Reference-Format style,
%% or BibLaTeX with the acmnumeric or acmauthoryear sytles, that include
%% support for advanced citation of software artefact from the
%% biblatex-software package, also separately available on CTAN.
%%
%% Look at the sample-*-biblatex.tex files for templates showcasing
%% the biblatex styles.
%%

%%
%% The majority of ACM publications use numbered citations and
%% references.  The command \citestyle{authoryear} switches to the
%% "author year" style.
%%
%% If you are preparing content for an event
%% sponsored by ACM SIGGRAPH, you must use the "author year" style of
%% citations and references.
%% Uncommenting
%% the next command will enable that style.
%%\citestyle{acmauthoryear}


%%
%% end of the preamble, start of the body of the document source.
\begin{document}
\newcommand{\TODO}[1]{\textcolor{red}{\textbf{TODO:} #1}}
\newcommand{\SystemName}[1]{SURGE} % we can change this for submission in case it's also in your thesis.
%%
%% The "title" command has an optional parameter,
%% allowing the author to define a "short title" to be used in page headers.

%% TODO: title
\title{\SystemName{}: Updating billion-scale vectors on distributed indexes}

%%
%% The "author" command and its associated commands are used to define
%% the authors and their affiliations.
%% Of note is the shared affiliation of the first two authors, and the
%% "authornote" and "authornotemark" commands
%% used to denote shared contribution to the research.
\author{Lars Th{\o}rv{\"a}ld}
\affiliation{%
  \institution{The Th{\o}rv{\"a}ld Group}
  \city{Hekla}
  \country{Iceland}}
\email{larst@affiliation.org}

\author{Valerie B\'eranger}
\affiliation{%
  \institution{Inria Paris-Rocquencourt}
  \city{Rocquencourt}
  \country{France}
}

\author{Aparna Patel}
\affiliation{%
 \institution{Rajiv Gandhi University}
 \city{Doimukh}
 \state{Arunachal Pradesh}
 \country{India}}

\author{Huifen Chan}
\affiliation{%
  \institution{Tsinghua University}
  \city{Haidian Qu}
  \state{Beijing Shi}
  \country{China}}

%%
%% By default, the full list of authors will be used in the page
%% headers. Often, this list is too long, and will overlap
%% other information printed in the page headers. This command allows
%% the author to define a more concise list
%% of authors' names for this purpose.
\renewcommand{\shortauthors}{Trovato et al.}

%%
%% The abstract is a short summary of the work to be presented in the
%% article.
\begin{abstract}
Graph-based approximate nearest neighbor (ANN) indexes such as HNSW
deliver the best recall--latency trade-offs for high-dimensional vector
search, but their reliance on fine-grained local connectivity makes them
difficult to distribute across machines and to update online. Existing
distributed graph-based systems route queries through a partitioning and
routing layer that is computed once from a static snapshot of the data and
cannot absorb insertions or deletions without an expensive full reindex.
We present \SystemName{}, a distributed ANN system that retains the
recall--latency benefits of HNSW while supporting continuous updates. \SystemName{}
combines (i) a neighborhood-preserving partitioning built over the base
layer of a routing HNSW that is maintained online through incremental
centroid updates, (ii) an adaptive maintenance policy that triggers a
coordinated, low-movement re-partitioning only when centroid drift exceeds
a configurable threshold, and (iii) three routing modes, including a
recall-targeted mode that decides which partitions to probe before
dispatch and therefore preserves full cross-worker parallelism. To
evaluate the system, we extend the NeurIPS~2023 Big-ANN streaming
generator to produce billion-scale runbooks with clustered, uniform, and
distribution-shifting update patterns. On SIFT and MSTuring workloads
across a ten-worker cluster, \SystemName{} sustains high recall@10 under
clustered and shifting updates where a static routing layer degrades by
up to \TODO{measure: max recall drop without maintenance, e.g. 30}\,percentage points, matches the static
throughput of the state-of-the-art distributed system GP-ANN within
\TODO{measure: \% throughput gap at equal recall} while additionally
supporting updates, and recovers nearly all of the recall of a full
per-step reindex at a fraction of its cost.
\end{abstract}

% %%
% %% The code below is generated by the tool at http://dl.acm.org/ccs.cfm.
% %% Please copy and paste the code instead of the example below.
% %%
% \begin{CCSXML}
% <ccs2012>
%  <concept>
%   <concept_id>00000000.0000000.0000000</concept_id>
%   <concept_desc>Do Not Use This Code, Generate the Correct Terms for Your Paper</concept_desc>
%   <concept_significance>500</concept_significance>
%  </concept>
%  <concept>
%   <concept_id>00000000.00000000.00000000</concept_id>
%   <concept_desc>Do Not Use This Code, Generate the Correct Terms for Your Paper</concept_desc>
%   <concept_significance>300</concept_significance>
%  </concept>
%  <concept>
%   <concept_id>00000000.00000000.00000000</concept_id>
%   <concept_desc>Do Not Use This Code, Generate the Correct Terms for Your Paper</concept_desc>
%   <concept_significance>100</concept_significance>
%  </concept>
%  <concept>
%   <concept_id>00000000.00000000.00000000</concept_id>
%   <concept_desc>Do Not Use This Code, Generate the Correct Terms for Your Paper</concept_desc>
%   <concept_significance>100</concept_significance>
%  </concept>
% </ccs2012>
% \end{CCSXML}

% \ccsdesc[500]{Do Not Use This Code~Generate the Correct Terms for Your Paper}
% \ccsdesc[300]{Do Not Use This Code~Generate the Correct Terms for Your Paper}
% \ccsdesc{Do Not Use This Code~Generate the Correct Terms for Your Paper}
% \ccsdesc[100]{Do Not Use This Code~Generate the Correct Terms for Your Paper}

% %%
% %% Keywords. The author(s) should pick words that accurately describe
% %% the work being presented. Separate the keywords with commas.
% \keywords{Do, Not, Use, This, Code, Put, the, Correct, Terms, for,
%   Your, Paper}
% %% A "teaser" image appears between the author and affiliationwher
% %% information and the body of the document, and typically spans the
% %% page.

% \received{20 February 2007}
% \received[revised]{12 March 2009}
% \received[accepted]{5 June 2009}

% %%
% %% This command processes the author and affiliation and title
% %% information and builds the first part of the formatted document.
\maketitle

\section{Introduction}
% Information is increasingly being stored as high dimensional vector embeddings in domains such as retrieval-augmented generation (RAG) for machine learning, recommendation systems, computer vision, and search \TODO{cite}. 
% These embeddings capture semantic relationships between data points, enabling similarity search via k-nearest neighbor (kNN) queries.

Modern applications increasingly rely on high-dimensional vector embeddings to represent data in domains such as retrieval-augmented generation (RAG) \cite{lewis2020retrieval, gao2023retrieval}, recommendation systems \cite{rajput2023recommender, okura2017embedding}, and large-scale search \cite{johnson2019billion, gordo2016deep}. These embeddings enable similarity search through $k$-nearest neighbor (kNN) queries, which are often latency-critical operations in production systems \cite{jayaram2019diskann, singh2021freshdiskann}.

At scale, where datasets contain billions of vectors, exact kNN search becomes computationally expensive. Instead, practical systems rely on approximate nearest neighbor (ANN) algorithms, which trade a small amount of accuracy for significant reductions in search latency. Existing ANN algorithms fall into two general categories: graph-based indexes \cite{malkov2018efficient, jayaram2019diskann, wang2021comprehensive} and partition-based indexes \cite{jegou2011searching, baranchuk2018revisiting}, each offering distinct trade-offs in performance and adaptability. As datasets scale beyond the capacity of the largest single node, scale-out, distributed systems are required~\cite{deng2019pyramid, chen2021spann, xu2023spfresh, mohoney2025quake, gottensuren2025unleashing}. Despite its growing practical importance, the distributed setting remains comparatively understudied: most work on ANN targets a single machine, and the few distributed graph-based systems assume that indexes store static data. % (YC, resolved) added framing that distributed dynamic ANN is important but understudied.

Partition-based indexes organize vectors into clusters and answer queries by probing a subset of partitions. Partitioned indexes are attractive due to their simplicity and flexibility. Distributed vector search systems commonly adopt them because they are straightforward to shard across machines and support efficient updates without requiring structural changes \cite{chen2021spann, xu2023spfresh, mohoney2025quake}. However, this design has fundamental limitations: the number of partitions to create and probe must be carefully decided before the index is built to achieve the desired performance. Query performance is bottlenecked by the need to perform full scans of multiple partitions, which leads to higher search latency compared to graph indexes \cite{wang2021comprehensive, aumuller2020ann, li2020annsurvey}.

Graph-based indexes build connections between nodes that enable efficient greedy traversal to quickly converge to high-quality nearest neighbors~\cite{malkov2018efficient, fu2019nsg, jayaram2019diskann}. As a result, graph-based methods achieve superior recall-latency trade-offs~\cite{li2020annsurvey, wang2021comprehensive, aumuller2020ann}. Despite the performance advantages, their reliance on fine-grained connectivity makes them difficult to deploy in distributed environments. Their performance depends on connectivity between close neighbors, and partitioning can degrade search quality by breaking graph links and disrupting local neighborhood structure~\cite{gottensuren2025unleashing, deng2019pyramid}.

Existing approaches attempt to bridge this gap: they combine dataset partitioning with graph-based search by introducing a routing layer over multiple local graph indexes \cite{gottensuren2025unleashing, deng2019pyramid}. Although these approaches improve scalability, they support only static datasets---that is, they provide no mechanism for inserting or deleting vectors after the index is built. In particular, both the indexes in~\cite{gottensuren2025unleashing} and Pyramid~\cite{deng2019pyramid} build static partitioning and routing schemes, which are unable to handle updates to data stored in the index. Such indexes would require expensive rebuilding to capture shifts in data distribution caused by update operations~\cite{singh2021freshdiskann, xu2023spfresh, mohoney2025quake}. Supporting updates in distributed graph-based systems remains challenging due to the need to preserve both partition quality and connectivity while minimizing data movement and expensive re-indexing operations. These challenges limit the deployment of graph-based indexes in high-speed, large-scale, dynamic systems.
% (YC, resolved) clarified that "static data" means no insert/delete support.

In this paper, we address these limitations by designing \SystemName{}, a distributed ANN system that retains the performance benefits of graph-based indexes while supporting efficient updates. We focus on Hierarchical Navigable Small World (HNSW) graphs and introduce a framework that enables their deployment in distributed settings without sacrificing recall or latency.

Our primary contributions are:
\begin{itemize}
    \item {\bf Contribution 1: Adaptive and neighborhood-preserving distributed partitioning.} We design a strategy that assigns vectors to partitions such that nearest neighbors are co-located, enabling high-recall graph-based search within each partition. Unlike prior approaches that compute partitions offline from a static snapshot~\cite{gottensuren2025unleashing, deng2019pyramid, chen2021spann}, our partitions remain representative as data is inserted and deleted, tracking shifts in the underlying distribution without requiring a full reindex.

    \item {\bf Contribution 2: Adaptive maintenance policy.} We introduce a maintenance policy that monitors partition quality over time and triggers updates to the routing layer and local graphs when degradation exceeds a configurable threshold.

    \item {\bf Contribution 3: Adaptive routing.} We develop a routing algorithm that directs each query to the partitions most likely to contain its nearest neighbors under the current data distribution. We support three routing modes: a fixed-fanout mode in which the number of candidate centroids used for routing is set by the user, a fixed partition mode which specifies the exact number of partitions to visit, and a recall-targeted mode in which the router dynamically determines how many partitions to probe in order to meet a specified recall target.

    \item {\bf Contribution 4: Large scale, dynamic evaluation with a dynamic benchmark generator.} We construct benchmarking workloads using existing billion-scale datasets, including SIFT and MSTuring~\cite{simhadri2022results, simhadri2024results}, to evaluate system performance under realistic update and query patterns. We extend the NeurIPS~2023 Big-ANN streaming workload generator to emit billion-scale runbooks with clustered, uniform, and distribution-shifting update patterns, and release it alongside the system so that distributed dynamic ANN systems can be evaluated under controlled distribution shift. We compare against both distributed static systems and single-machine dynamic systems \TODO{may not include Quake at the end}.

    \item {\bf Contribution 5: \yc{build and release a system? lets talk about this} }
\end{itemize}

We evaluate \SystemName{} against GP-ANN~\cite{gottensuren2025unleashing} (distributed, static), Quake~\cite{mohoney2025quake} (single-node, dynamic), and a distributed PostgreSQL deployment with the \texttt{pgvector} and Citus extensions~\cite{postgresql, pgvector, cubukcu2021citus} (distributed, dynamic). \SystemName{} sustains high recall@10 on clustered and shifting workloads where static routing degrades by up to \TODO{measure: max recall drop without maintenance}\,percentage points, delivers \TODO{measure: X}$\times$ the query throughput of distributed PostgreSQL at \TODO{measure: Y}\,\% recall, and matches GP-ANN's static throughput within \TODO{measure: Z}\,\% at equal recall while additionally supporting online updates.

% We compare with several existing systems to test our partitioning methods, maintenance policies, and routing algorithms.




% Graph indexes construct proximity graphs where each node is connected to its approximate neighbors. Partitioned indexes use a clustering algorithm to sort vectors into partitions. Queries are processed by scanning a subset of the partitions. Partitioned indexes are attractive due to their simplicity and flexibility. Distributed vector search methods typically use partitioned indexes because they are easy to shard across machines. Additionally, these indexes allow for simple update operations since they do not require updating connections within a graph structure. However, graph indexes offer the best throughput versus accuracy performance. \TODO{transition} Despite these advantages, graph-based methods present challenges in distributed environments. Their performance depends on carefully maintained connectivity, and naive partitioning can degrade search quality by breaking graph links across machines. Furthermore, supporting dynamic updates--such as insertions and deletions--introduces additional complexity compared to partitioned approaches.
% For a query, Faiss HNSW can achieve 10x performance over Faiss IVF index



\section{Background \& Motivation}

\subsection{Approximate Nearest Neighbor Search}

    Vector similarity search is the task of retrieving the $k$ most similar vectors to a query $q$ from a dataset $\mathbf{X} = {x_1,...,x_n} \subset \mathbb{R}^d$ under a distance metric such as Euclidean distance or inner product. Exact $k$-nearest neighbor search is computationally prohibitive at scale, motivating approximate nearest neighbor (ANN) algorithms that trade a small loss in accuracy for significant gains in query throughput and latency~\cite{malkov2018efficient, jegou2011searching, li2020annsurvey, aumuller2020ann}.

    Among ANN algorithms, graph-based indexes such as Hierarchical Navigable Small Worlds (HNSW)~\cite{malkov2018efficient} have emerged as performance leaders, especially for high-dimension embeddings. By building a proximity graph and answering queries via greedy traversal, these methods achieve recall-latency trade offs that partition-based approaches have difficulty matching~\cite{malkov2018efficient, jayaram2019diskann, li2020annsurvey, aumuller2020ann}. The key property of graph indexes is locality: vectors that are close in the embedding space are connected in the graph, so greedy search converges quickly with few distance computations. Partition-based indexes such as IVF~\cite{sivec2003ivf,jegou2011searching,douze2025faisslibrary} probe a subset of clusters at query time. They are simple, memory-efficient, and naturally amenable to sharding and updates, which makes them the common choice for distributed and dynamic deployments. Their main cost is that each probed cluster is scanned exhaustively, so reaching the recall of a graph index typically requires many more distance computations. \SystemName{} aims to combine the update-friendliness of partition-based systems with the search efficiency of graph indexes. % (YC, resolved) softened the framing of partition-based methods and stated the complementary goal rather than a head-to-head dismissal.

\subsection{Distributed ANN}
    % (YC, resolved) added coverage of non-HNSW distributed methods (SPANN/SPFresh, partition-replication systems) before focusing on the graph-based two-level architecture.
    Distributed ANN systems fall broadly into two families. Partition-based systems such as SPANN~\cite{chen2021spann} and its dynamic successor SPFresh~\cite{xu2023spfresh} shard vectors into many posting lists, replicate boundary points, and answer queries by scanning a chosen number of partitions; they update cheaply but inherit the exhaustive-scan cost of IVF. Graph-based systems instead build a routing layer over the dataset and are the focus of this work, because they retain the superior recall--latency trade-off of single-node graph indexes. Scaling graph-based indexes to datasets that exceed single-machine memory requires distributing the index across a cluster. Naive partitioning---randomly assigning vectors across nodes and broadcasting every query---preserves recall but eliminates throughput gains, since every node must perform a search for every query. The fundamental difficulty is that graph-based search depends on fine-grained local connectivity: the greedy traversal works because nearby vectors are linked, and a random partition severs these links at partition boundaries, degrading both recall and routing efficiency. Any practical distributed graph-based system must therefore solve two problems: how to assign vectors to partitions such that nearest neighbors are co-located, and how to route each query to only the partitions likely to contain its answer.
    
    % Any practical distributed graph-based system must therefore solve a partitioning problem: how to assign vectors to shards such that nearest neighbors are co-located, and how to route each query to only the shards likely to contain its answer.

    Prior work~\cite{deng2019pyramid, gottensuren2025unleashing} has addressed this problem with a two-level architecture.
    % The principled solution to this problem, introduced by Pyramid \TODO{cite} and developed further by GP-ANN \TODO{cite}, is a two-level architecture. 
    Rather than routing queries based on random or round-robin assignment, a small routing index is built over a set of representative centroids that captures the global structure of the dataset. At query time, the router searches this index to identify which partitions are relevant to the query, and forwards the query only to those nodes. Since the routing index is small, this step is fast and because the partitioning was constructed to co-locate nearby vectors, the selected partitions are likely to contain the true nearest neighbors.

    % The Hierarchical Navigable Small World (HNSW) graph \TODO{cite} is the dominant graph-based ANN algorithm, organizing vectors in a multi-layer structure that enables greedy search in sub-linear time \TODO{cite}. While HNSW achieves excellent single-node performance, scaling it to datasets that exceed single-machine memory requires distributing the index across nodes \TODO{cite}. Naive sharding -- randomly assigning vectors across nodes and broadcasting every query -- preserves recall but eliminates throughput gains, since every node must perform a search for every query.

    % Pyramid \TODO{cite} introduced a two-level architecture for distributed HNSW that avoids this bottleneck. A small \textit{meta-HNSW} graph is built over centroids representative of the data distribution and serves as a routing layer. The base layer of the meta-HNSW graph, which contains all centroids, is further partitioned to match the number of available nodes. Each node is assigned a partition and builds a local HNSW graph over its assigned vectors. At query time, the routing node searches the meta-HNSW to find the K nearest centroids and then routes the query to the corresponding worker nodes. This design achieves high throughput by contacting only a small fraction of workers per query while preserving recall through neighborhood-preserving partitioning. \TODO{architecture fig?}

    % GP-ANN \TODO{cite} extended this two-level approach with an alternative routing and partitioning scheme. \TODO{explain gp-ann briefly}. 
    
    % Both Pyramid and GP-ANN construct static indexes: the partitioning and routing scheme is computed once from a fixed snapshot of the data and cannot adapt to subsequent insertions or deletions without a full and expensive reindex. This limitation is the primary gap our work addresses.

    The quality of this routing scheme depends critically on how well the centroid positions and partition boundaries reflect the actual distribution of vectors in the system. When the dataset is static, this is a one-time cost: the routing index is built once and remains valid indefinitely. Both Pyramid and GP-ANN operate under exactly this assumption, constructing their routing layers offline from a fixed snapshot of the data and providing no mechanism for updating them as vectors are inserted or deleted. Any change to the dataset requires a full reindex. This is the central limitation our work addresses.

\subsection{Dynamic ANN Systems}
    Supporting efficient updates in ANN indexes is difficult even on a single node. The core challenge for graph-based indexes is that deletions are structurally destructive: removing a node severs edges that may be on navigating paths to other nodes, degrading the reachability of the graph and causing recall to drop over time~\cite{singh2021freshdiskann, yu2025topology, xiao2024enhancing, xu2025place}.
    Alternatively, shadow deletes or tombstones can be used to mark deleted nodes as inactive without modifying the graph structure~\cite{malkov2018efficient, singh2021freshdiskann}. However, these deleted vectors remain in memory and on navigating paths indefinitely, increasing traversal cost and memory footprint without bound as deletions accumulate.
    % Simple strategies such as shadow deletes and tombstones mark nodes as inactive without modifying the graph, but allow deleted vectors to consume memory and participate in graph traversal indefinitely, causing unbounded degradation under high delete rates. 
    This issue has largely been addressed on single node graph index deployments by introducing graph update rules that preserve the navigability of the index across a stream of insertions and deletions~\cite{yu2025topology, xu2025place, xiao2024enhancing}.

    In the distributed setting, the update problem acquires an additional dimension. Beyond maintaining local graph quality within each partition, updates shift the underlying data distribution in ways that degrade the routing layer. As new vectors are inserted into regions of the space that were previously sparse, the centroids of the routing index become unrepresentative, and partition boundaries no longer reflect where nearby vectors actually reside. Inserted vectors may then repeatedly hit the same partitions, causing unbalanced load across nodes while queries are routed to the wrong partitions, causing recall to degrade. This routing degradation is distinct from the local graph degradation studied by prior single-node solutions, and it cannot be addressed by local graph maintenance alone. Quake \cite{mohoney2025quake} recognizes that index structure must adapt to shifting data distributions and introduces a cost-model-driven approach to dynamic partitioned index maintenance on a single node, but does not extend to the distributed setting where data must be repartitioned and rebalanced across machines. No existing system jointly addresses routing maintenance and local update support in a distributed graph-based index.

    % Supporting efficient updates in ANN indexes is an active area of research, with some existing solutions already addressing the issue of updates and shifting workloads.

    % FreshDiskANN \TODO{cite} extends DiskANN \TODO{cite} to support real time updates by maintaining an append-only structure for new vectors on disk and periodically merging them into the main graph. It demonstrates strong performance on the SIFT dataset, although SIFT data has a uniform distribution and does not properly stress-test the routing layer: new insertions fall into the existing distribution. \TODO{VERIFY CLAIMS}

    % Quake \TODO{cite} introduces adpative single-node indexing with a cost model that dynamically selects execution parameters to meet recall and latency targets. It relies on an IVF-style index which can accept update operations. Under skewed and evolving workloads, Quake has mechanisms to perform index maintanence to dynamically adjust the number of partitions and reorganize data within them. However, Quake currently operates as a single node system, relies on splitting data into thousands to tens of thousands of partitions, and performs brute-force search within each partition. This approach does not take advantage of the latency benefits graph indexes provide and limits the deployment of Quake on large datasets.

\subsection{Dynamic Workload Benchmarks}

    Until recently, ANN evaluation assumed static datasets: an index is built once and repeatedly queried~\cite{simhadri2022results,simhadri2024results,jegou2011searching,yandex2016deep}. This assumption is increasingly unrealistic as production workloads insert and delete vectors continuously~\cite{singh2021freshdiskann, xu2023spfresh, mohoney2025quake}. The NeurIPS~'23 Big-ANN Challenge~\cite{simhadri2024results} introduced the first standardized benchmark for evaluating ANN systems under streaming workloads which interleave insertions, deletions, and queries. The benchmark contains 'runbooks' which define sequences of searches, insertions, and deletions. Updates are performed in clustered batches that progressively shift the data distribution. This is specifically useful to test the adaptability of routing layers in distributed indexes: recall can degrade as insertions that fall far from existing centroids are systematically misrouted. We use this benchmark to generate runbooks of both clustered and uniform updates for larger datasets to evaluate our own system. We use similar generation scripts to also generate runbooks with even larger distribution shift, which are not currently available in the Big-ANN generator, to further stress-test the routing layer and maintenance policies of our system.
    
    % , as well as the methods for generating it \TODO{generate larger datasets and runbooks in the same way}, to evaluate our own system. 

\subsection{Problem Motivation}
\label{sec:motivation}
    % (YC, resolved) kept Problem Motivation as a distinct subsection that ties together the three failure modes below; if space allows, promote Figure~\ref{fig:motivation} to page 1 and reference it from the Introduction. The subsection can be lifted out of Section 2 into its own section without changes if the editors prefer.
    The two-level distributed architecture above is efficient only when the routing layer reflects the current data distribution. Under updates this assumption breaks in three distinct ways, none of which is repaired by local graph maintenance alone.

    First, on a single node, a graph index whose only update mechanism is shadow deletion accumulates inactive nodes on navigating paths, so traversal cost grows and recall drifts down as deletions accumulate~\cite{singh2021freshdiskann, yu2025topology}. This motivates per-partition graph maintenance, but it is not the dominant effect at scale.

    Second, in the distributed setting, clustered insertions move the true data distribution away from the centroids that the routing layer was built on. Because the router selects partitions purely from centroid positions, queries are increasingly dispatched to the wrong partitions and recall falls. This routing degradation can only be addressed by a full index maintenance operation.
    % ---\emph{even when every probed partition is searched exhaustively and the local graphs are perfect}

    Third, routing a stream of clustered insertions to their nearest existing partition without re-partitioning concentrates new vectors on a few workers. The resulting load imbalance erodes the throughput advantage of distribution, because the most-loaded worker bounds end-to-end latency.

    \begin{figure}
        \centering
        % \includegraphics[width=\linewidth]{figures/motivation.pdf} % TODO: generate and uncomment
        \fbox{\parbox[c][3.2cm][c]{0.92\linewidth}{\centering\textcolor{red}{\textbf{Placeholder: motivating figure.}}\\ \footnotesize Recall, QPS, and partition imbalance vs.\ runbook step for single-node HNSW, distributed static routing, and naive distributed inserts on a clustered workload.}}
        \caption{\textbf{Motivation.} On a clustered update workload, single-node HNSW with shadow deletes loses recall as deleted nodes accumulate; a distributed index with a \emph{static} routing layer loses recall because insertions shift the data away from the fixed centroids, \emph{even though each probed partition is searched exhaustively}; and naive insertion without re-partitioning drives up partition imbalance and depresses throughput. Local graph maintenance addresses only the first failure mode, motivating \SystemName{}'s online routing maintenance.}
        \label{fig:motivation}
    \end{figure}

    \TODO{Figure~\ref{fig:motivation}: recall, QPS, and partition imbalance vs.\ runbook step on a clustered workload for (1) single-node HNSW with shadow deletes, (2) distributed system with static routing, and (3) distributed system with naive inserts and no re-partitioning. The figure should show that static routing degrades recall on clustered updates even when local search is exhaustive, and that naive inserts inflate imbalance/lower throughput. Optionally add a 2-D illustration contrasting a uniform dataset with a clustered/shifting one and the resulting partitions.}

    \subsubsection{Evaluation Metrics}
    \label{sec:eval-metrics}

    We introduce \textit{theoretical recall} as a metric that isolates routing quality from local search quality. Formally, given a query $q$ with ground-truth nearest neighbors $N^*(q)$, and a routing decision that selects a set of partitions $\mathcal{P}(q)$, we define:
    
    \begin{equation}
    \mathrm{TheoreticalRecall}@k(q) \;=\; \frac{\bigl|\,\{x \in N^*(q) \;:\; x \in \bigcup_{P \in \mathcal{P}(q)} P\}\,\bigr|}{k},
    \label{eq:theoretical-recall}
    \end{equation}
    where $N^*(q)$ is the set of $q$'s true $k$ nearest neighbors and $\mathcal{P}(q)$ is the set of partitions the router selects for $q$.
    Theoretical recall measures the fraction of true nearest neighbors that reside in the partitions selected by the router, under the assumption that each partition is searched exhaustively. It is an upper bound on the recall any local search algorithm can achieve given the routing decision: a system cannot find a nearest neighbor that was routed to the wrong partition. This metric allows us to evaluate the routing pipeline---the choice of centroid representation, routing graph construction, and partitioning algorithm---independently of local search.

% , and to identify precisely when and why routing maintenance is necessary.
% \begin{itemize}
%     \item ANN fundamentals
%     \item Distributed HNSW
%         \begin{itemize}
%             \item Pyramid architecture
%             \item GP-ANN routing \& partitioning
%         \end{itemize}
%     \item Dynamic ANN systems
%         \begin{itemize}
%             \item FreshDiskANN
%             \item Quake
%         \end{itemize}
%     \item Streaming benchmarks - neurips '23
%     \item problem formulation/motivation
%         \begin{itemize}
%             \item run single node hnsw on streaming benchmark to demonstrate degradation of performance
%             \item show that no change in routing works but leads to imbalance? \TODO{NOW}
%             \item define (theoretical) recall
%         \end{itemize}
% \end{itemize}

\section{\SystemName{}}
    \SystemName{} adopts a two-tier  coordinator-worker architecture. A single coordinator node maintains a routing index built over cluster centroids representative of the global data distribution. The centroids are indexed in a HNSW graph for fast routing. Each of the worker nodes maintains a local HNSW graph over the vectors assigned to its partition. This routing architecture is the same as that used by Pyramid \cite{deng2019pyramid} and GP-ANN \cite{gottensuren2025unleashing}, which find that an HNSW graph as the router works well. The coordinator receives all queries and update requests, routes them to the appropriate worker(s), and orchestrates index maintenance. Workers serve local searches, apply updates, and participate in coordinated index maintenance operations.

    \begin{figure}
        \centering
        % \includegraphics[width=\linewidth]{figures/architecture.pdf} % TODO: generate and uncomment
        \fbox{\parbox[c][3.6cm][c]{0.92\linewidth}{\centering\textcolor{red}{\textbf{Placeholder: architecture diagram.}}\\ \footnotesize Coordinator holding the routing HNSW $G_m$, live centroids $C$, and routing table $\rho$; $w$ workers each holding a local HNSW $G_i$. Arrows show query routing, update forwarding, and maintenance-time vector exchange.}}
        \caption{\textbf{\SystemName{} architecture.} A single coordinator maintains a routing HNSW $G_m$ over $m$ representative centroids, a live copy of each centroid position, and a routing table $\rho$ mapping centroids to workers. Each of the $w$ workers stores a local HNSW $G_i$ over its assigned partition. The coordinator routes every query and update to the relevant worker(s) and orchestrates maintenance; workers serve local searches, apply updates, and exchange vectors during re-partitioning.}
        \label{fig:architecture}
    \end{figure}

    Figure~\ref{fig:architecture} illustrates the overall architecture.

\subsection{Routing Layer Construction}
    The routing layer is constructed once during index initialization and subsequently maintained online as vectors are inserted and deleted. It consists of a routing HNSW graph over a set of representative centroids and a routing table mapping each centroid to its owning worker.

\subsubsection{Graph Construction}
    The coordinator begins by drawing a random sample $X' \subset X$ of size $n' \ll n$ from the dataset. A $k$-means clustering with $m$ centers is run on $X'$ to produce a set of centroids $C = \{c_1, \ldots, c_m\}$ representative of the underlying data distribution. 
    % Each centroid $c_i$ is assigned a weight equal to the number of sample points it attracts, so that subsequent partitioning can account for load balance.

    The coordinator then constructs an HNSW graph $G_m$ over the $m$ centroids using the standard HNSW construction algorithm~\cite{malkov2018efficient}. Because $|C| \ll n$, building $G_m$ is inexpensive and the resulting graph is small enough to be used for fast query routing. Pyramid finds that $m = 10^4$ centroids is sufficient for high routing accuracy across the datasets we consider; larger routing HNSWs increase construction and search latency without meaningfully improving recall~\cite{deng2019pyramid}.

\subsubsection{Graph Partitioning}
    \label{sec:partitioning}
    To assign centroids to workers, we extract the base layer of $G_m$---a proximity graph over all $m$ centroids---and partition it into $w$ balanced parts, one per worker. Partitioning the graph containing nearest-neighbor connections rather than the centroids directly ensures that centroids assigned to the same worker are close in the embedding space. This ensures that nearby vectors are placed in the same partition and attempts to preserve the local neighborhood structure that would be present in a single node system.
    
    We use the KaHIP graph partitioner~\cite{sandersschulz2013kahip}, which minimizes the number of edges crossing partition boundaries subject to a balance constraint on total vertex weight. Optionally, the weight of each centroid can be set to reflect the number of dataset vectors that were routed through it to the workers for insertion, so the partitioning simultaneously minimizes cross-partition nearest-neighbor placement and balances the per-worker vector count. 
    % For initial index construction we use KaHIP's \texttt{STRONG} mode; for online repartitioning during maintenance (Section~\ref{sec:maintenance}) we use \texttt{FAST} to reduce overhead.
    
    Once partitioning is complete and each centroid has been assigned a partition, the coordinator builds a global routing table $\rho: C \rightarrow \{1, \ldots, w\}$ mapping each centroid to its assigned worker. Both $G_m$ and $\rho$ are stored on the coordinator, which uses the two to route every batch of queries.
    
    To complete the initial build phase, the coordinator then distributes vectors to workers: for each $x \in X$, the coordinator searches $G_m$ to find $x$'s nearest centroid $c^*$ and forwards $x$ to worker $\rho(c^*)$. Each worker independently builds a local HNSW $G_i$ over its assigned vectors.
    
\subsection{Query Routing}
    Currently, \SystemName{} supports three methods of executing search queries across workers: 

\subsubsection{Static Routing K}
    Pyramid~\cite{deng2019pyramid} utilizes a static routing factor $K$ to route queries to the closest partitions - this $K$ factor is used for the top level routing graph only. Upon receiving a $k$-nearest-neighbor query $q$, the coordinator searches the routing HNSW $G_m$ to retrieve the top-$K$ nearest centroids $\mathcal{N}_K(q) = \{c_{i_1}, \ldots, c_{i_K}\}$. The set of target workers, $\omega(q)$,  is then determined by applying the routing table:
    \[
      \omega(q) = \left\{ \rho(c) \mid c \in \mathcal{N}_K(q) \right\}.
    \]
    The coordinator dispatches $q$ to each worker in $\omega(q)$, which independently search their local HNSWs and return their local top-$k$ candidates. The coordinator merges these result sets and returns the global top-$k$.

    The routing parameter $K$ controls the recall-throughput trade-off with a larger granularity than simply declaring a static number of partitions to search. With $K = 1$, each query touches a single partition. Increasing $K$ broadens the search to additional partitions and improves recall, at the cost of increased per-query work. Because the partitioning co-locates proximate centroids, the $K$ nearest centroids in $G_m$ often fall within the same or a small number of partitions; accordingly, $|\omega(q)| \leq K$ and is typically well below the total partition count even for large $K$ \cite{deng2019pyramid}.

\subsubsection{Enforced Nprobe}
    \label{sec:enforced-nprobe}
    Many distributed and partition-based ANN systems probe a fixed number of partitions $N$ per query, where $N$ is a tunable knob that trades recall for latency~\cite{gottensuren2025unleashing, mohoney2025quake, jegou2011searching, chen2021spann}. \SystemName{} supports this routing mode for compatibility with such systems and to enable apples-to-apples recall comparisons. Given a target $N$, the coordinator searches the routing HNSW $G_m$ for the closest centroids, maps them to partitions using the routing table $\rho$, and checks whether the results cover at least $N$ distinct partitions. Because the partitioning ensures proximate centroids fall in the same partitions (Section~\ref{sec:partitioning}), the top-ranked centroids often fall within fewer than $N$ partitions. In that case the coordinator must re-query $G_m$ with a larger $K$ until the centroids returned map to at least $N$ partitions.
    
    Repeating the routing search is wasteful: each retry rediscovers centroids the previous search already returned, and the overhead grows with $N$. Lu et al.~\cite{lu2026depth} address this with a \emph{resumable} HNSW search that retains the candidate queue between query calls and continues traversal from where it left off, paying only the marginal cost of expanding the frontier. \texttt{hnswlib} does not have this functionality, and although adding it is straightforward, we leave it as future work. Enforced-nprobe routing in \SystemName{} is therefore used only for direct recall comparisons against baselines whose only routing mode is a fixed nprobe; other routing modes are used to measure throughput performance. 

% \subsubsection{Dynamic Routing}
%     In addition to the static modes described above, \SystemName{} supports a routing mode which aims to reach a particular recall target, so the number of partitions searches is determined adaptively per query. Rather than routing to a fixed number of centroids or partitions, the coordinator expands the search frontier to more partitions until the routing decision is expected to satisfy a user-specified recall target.

%     Quake \cite{mohoney2025quake} implements Adaptive Partition Scanning (APS) using a geometric model to estimate the probability that each partition contains one of the query's k nearest neighbors. However, APS requires that partitions be searched sequentially, which in our case would limit throughput and lead to higher search latencies. We instead implement a simpler approximation to calculate the probability that a given partition contains the nearest neighbors of a query from the centers returned by the routing HNSW. 

%     \TODO{describe the specific adaptive algorithm}

\subsubsection{Dynamic Routing}
    \label{sec:dynamic-routing}
    In addition to the static modes above, \SystemName{} supports a routing mode that targets a user-specified recall target $\hat{r}$ and adaptively chooses how many and which partitions to probe per query. Rather than committing to a fixed routing $K$ or partition count, the coordinator orders partitions by an estimate of how likely each is to contain the query's nearest neighbors and stops once the cumulative estimate crosses $\hat{r}$.

    Quake~\cite{mohoney2025quake} implements a similar search algorithm with Adaptive Partition Scanning (APS), which scans partitions sequentially and updates a recall estimate from the intermediate top-$k$ results after searching each one, terminating once the estimate exceeds the target. Because the estimate depends on results from already-scanned partitions, APS must visit partitions sequentially. In a distributed setting, this would serialize work across workers and erase the parallelism that motivates partitioning in the first place.

    \SystemName{} instead decides which partitions to query \emph{before} dispatching, using only information already available on the coordinator: the centers returned by the routing HNSW and their distances to $q$. As a result, the scoring scheme we use use assumes that centroid positions are representative of the most current data distribution; when centroids are stale, the scores are miscalibrated because the router is unable to accurately compute scores.
    
    The coordinator retrieves the $K$ closest centers $\{c_1, \ldots, c_K\}$ to $q$ from $G_m$, where $K = 5 \cdot N_p$ and $N_p$ is the number of partitions. We found that setting $K$ to be a reasonable multiple of the total number of partitions ensures that even when many centers are co-assigned to the same partition, every partition has a chance of accumulating probability mass; a $K$ that is too small zeros out all but the nearest few partitions and makes high recall targets unattainable. Each center $c_r$ at rank $r$ and distance $d_r = d(q, c_r)$ contributes a weight to its owning partition $\rho(c_r)$:
%     \begin{equation}
%   \tilde{w}_r
%   \;=\;
%   \lvert\rho(c_r)\rvert
%   \cdot
%   \exp\!\left(-\alpha \,\frac{d_r}{d_0}\right),
%   \label{eq:center-weight}
% \end{equation}
    \[
    w(c_r) \;=\; \lvert\rho(c_r)\rvert \cdot \exp\!\left(-\frac{d_r}{d_0}\right),
    \]
    where $d_0$ is the distance to the single closest center and $\lvert\rho(c_r)\rvert$ is the number of vectors currently assigned to the partition of $c_r$. The unnormalized partition score aggregates the weights of all centers belonging to the same partition, $\tilde{p}(P) = \sum_{c \,:\, \rho(c) = P} w(c)$. And the normalized score $p(P) = \tilde{p}(P) / \sum_{P'} \tilde{p}(P')$ acts as a per-query estimate of the probability that $P$ contains one of the query's true nearest neighbors. The coordinator chooses partition to visit in descending order of $p(P)$, accumulating their scores, and stops as soon as the running sum reaches $\hat{r}$. All selected partitions are then queried in parallel
    
%     \begin{equation}
%   \tilde{p}(P)
%   \;=\;
%   \sum_{\substack{r=1 \\ \rho(c_r)=P}}^{K}
%   w(c_r),
%   \label{eq:partition-score}
% \end{equation}
    
    % \[
    % w(c_r) \;=\; \frac{1}{\bigl(d(q, c_i)\bigr)^{3}} \cdot \frac{1}{i},
    % \]
    % where $d(q, c_i)$ is the query-to-center distance and $i$ is the rank of $c_i$ in the routing-HNSW result (1-indexed). 
    
    % The unnormalized partition score is the sum of weights of its constituent centers, $\tilde{p}(P) = \sum_{c \,:\, \rho(c) = P} w(c)$, and the normalized score $p(P) = \tilde{p}(P) / \sum_{P'} \tilde{p}(P')$ acts as a per-query estimate of the probability that $P$ contains one of the query's true nearest neighbors. The coordinator visits partitions in descending order of $p(P)$, accumulating their scores, and stops as soon as the running sum reaches $\hat{r}$. All selected partitions are then queried in parallel.
    
    % The score combines two intuitions: a center close to $q$ is more likely to route to a partition with the true nearest neighbors of $q$, and a partition that owns several of the top-ranked centers offers stronger evidence than one that owns a single distant center. The specific functional form--cubic distance falloff combined with reciprocal-rank weighting--is empirical. We treated the distance exponent and the rank-weighting term as hyperparameters and selected them by sweeping each independently on several datasets; the combination above produced recall closest to the target across the $80$--$99\%$ range, which is the range relevant to most deployments. 

    The weight $w(c_r)$ combines two intuitions: the exponential $\exp(- d_r / d_0)$ is evaluated on the \emph{relative} distance $d_r / d_0$ rather than the raw distance. Normalizing by $d_0$ makes the distribution shape invariant to the absolute scale of distances in the embedding space and produces consistent partition-probability distributions across datasets with different characteristic inter-centroid distances. The factor $\lvert\rho(c_r)\rvert$ is a Bayesian prior on partition size: ignoring distances, a larger partition contains more vectors and is therefore more likely to contain the query's nearest neighbor. The coordinator maintains exact per-center counts as a byproduct of the incremental centroid updates described in Section~\ref{sec:update-handling}, so partition counts are readily available.

    A useful property of this score is that the neighbor count $k$ does not appear. Quake's geometric model depends on the radius to the current $k$-th nearest neighbor, which grows with $k$ and pulls additional partitions into the relevant set. \SystemName{}'s score instead estimates which partitions lie in the local neighborhood of $q$, not how many of $q$'s top-$k$ neighbors each partition holds. When $k$ is much smaller than the partition size, the partitions containing the top nearest neighbor are approximately the partitions containing the top hundred, so a single routing decision serves a wide range of $k$. We confirm this empirically: across $k \in \{1, 10, 100\}$, the same recall-target setting yields a theoretical recall within 5\% of the target on all three.
    
\subsection{Update Handling}
\label{sec:update-handling}
    \SystemName{} treats insertions and deletions as first-class operations. Unlike prior distributed graph-based systems that require a full offline reindex following any data change~\cite{deng2019pyramid, gottensuren2025unleashing}, our design supports continuous updates while preserving routing accuracy and local search quality.

\subsubsection{Update Routing}
    Each insertion or deletion modifies the underlying data distribution and therefore potentially shifts the optimal centroid positions. In a distributed system, a reindexing operation would require rebuilding the partitioning layer by collecting all the vectors from workers and recomputing the data distribution via $k$-means. To keep the routing layer current without requiring access to vectors on the workers or re-clustering, we maintain a live copy of each centroid on the coordinator and update it incrementally.

    \paragraph{Insertions.} When a vector $v$ is inserted, the coordinator searches $G_m$ to identify the current nearest centroid $c_i$ and updates its position as a moving average over the $n_i$ vectors currently assigned to it:
    \[
      \vec{c}_i \leftarrow \frac{\vec{c}_i \cdot n_i + \vec{v}}{n_i + 1}.
    \]
    The vector is then forwarded to the worker $\rho(c_i)$ for insertion into its local HNSW.
    
    \paragraph{Deletions.} Similarly, when a vector $v$ is deleted, the coordinator updates the centroid that the vector was initially routed through:
    \[
      \vec{c}_i \leftarrow \frac{\vec{c}_i \cdot n_i - \vec{v}}{n_i - 1}.
    \]
    The deletion is propagated to the owning worker as a shadow delete.
    
    Because both operations keep a copy of centroid positions that are consistent with the evolving data distribution, the degradation of the routing layer can be measured and corrected by maintenance operations when drift accumulates.
    
    % gracefully under moderate update rates and can be corrected by maintenance operations when drift accumulates.

\subsubsection{Maintenance}
    As updates accumulate, centroid drift can cause routing degradation and partition imbalance. \SystemName{} monitors partition quality continuously and triggers a coordinated maintenance operation when degradation exceeds a configurable threshold.
    
    After processing batches of updates, or at defined intervals, the coordinator reconstructs a candidate routing HNSW $G_m'$ from the updated centroids and re-partitions its base layer into $w$ parts. This recomputation is fast: the routing graphs consist of $10^4$ centroids \footnote{This value works for all datasets considered, as previously observed in Pyramid~\cite{deng2019pyramid}}, so building an HNSW graph on these centroids takes less than \TODO{measure: $G_m'$ build time, e.g. tens of ms} and the partitioning step takes \TODO{measure: KaHIP FAST partition time}. While this maintenance check takes place, the coordinator can continue serving queries.

    Because the graph partitioner assigns each new partition an arbitrary integer label $1,\ldots,w$ independently of the prior assignment, directly comparing old and new labels would overestimate the required data movement: a partition whose membership is nearly unchanged may simply receive a different label. % (resolved) defined "labels" as the partitioner's part identifiers.
    We therefore apply a maximum-weight bipartite matching between old and new partition labels, selecting the relabeling that maximizes centroid overlap between corresponding partitions to minimize data movement. Let $\phi$ denote the fraction of centroids whose partition assignment changes after optimal relabeling. If $\phi$ exceeds a given threshold $\tau$, all workers are notified and re-evaluate the partition assignment of each local vector and exchange vectors with other workers as needed. Once workers have their new working sets, they mark vectors that migrated to other workers as deleted in their local graphs, and insert vectors newly assigned to them. Using the \texttt{replace\_deleted} functionality in hnswlib ensures that graph size does not grow an unbounded amount. However, since the total number of vectors stored on each worker after maintenance can change significantly, after a certain percentage (we use 50\%) of slots in a local HNSW graph are inactive (ie have been marked deleted but not replaced) we force a full local graph rebuild, which builds an HNSW graph on only the current active vectors on a worker. Once the maintenance is complete, the coordinator and workers swap to the new routing and local indexes and subsequent queries use the newly constructed graphs.
    
    % the routing layer is swapped to the new meta-HNSW and the routing table is updated with the new partition assignment.

    Queries continue to be served from the old indexes throughout the rebuild. Updates arriving during a rebuild are buffered in an insert/delete log on the coordinator and applied after the rebuild completes.

    % \paragraph{Availability during rebuild.} The coordinator transitions through three states---\textsc{Idle}, \textsc{Rebuilding}, and \textsc{Draining}---to ensure uninterrupted query service. During \textsc{Rebuilding}, incoming updates are appended to the coordinator's log rather than forwarded to workers, keeping the old graphs stable for query processing. Once the new graphs are swapped in, the system enters \textsc{Draining} and flushes the buffered updates before accepting a new rebuild request.


\subsubsection{Updates to Local Graphs}
    In this section we discuss how local graphs handle the different update operations.

    \textit{Insertions.} Each worker's local graph natively supports incremental insertions. Upon receiving an insert request from the coordinator, a worker inserts the vector into its local graph using the standard HNSW insertion procedure. Concurrent inserts and queries are handled safely using the locking mechanisms provided by the underlying \texttt{hnswlib}~\cite{malkov2018efficient} implementation.

    \textit{Deletions.} Structural deletion from an HNSW graph requires modifying edges on all layers and risks degrading graph navigability~\cite{yu2025topology, xiao2024enhancing, xu2025place}. \SystemName{} instead employs shadow deletes with replacement: a deleted vector is marked inactive and excluded from search results, but its node and edges remain in the graph until either a vector is inserted and can replace it in the graph or the next full rebuild, at which point it is permanently removed. The delete and replace policy is adopted from \texttt{hnswlib}, which supports the option to replace deleted vectors with newly inserted ones. This approach preserves graph structure between rebuilds at the cost of a bounded increase in memory and traversal overhead proportional to the deletion rate.

\subsubsection{Performance}

    The time to complete a rebuild operation depends heavily on the amount of data movement that needs to occur. Additionally, there exist several options for local graph maintenance that have distinct trade offs. One option is for each worker to perform a full rebuild of the local graph with its newly reassigned vectors, however this option is the most time consuming and computationally expensive. Alternatively, local graphs can modify the existing graph structure and use shadow deletes to mark nodes that no longer belong to them as inactive while inserting newly assigned vectors. While this option is lightweight, it leads to a monotonic increase in the local graph sizes and if graphs become large enough, it negates the benefits of having a distributed system. A secondary option is to use the hnswlib-supported option to \texttt{replace\_deleted} vectors in the local graph with newly inserted vectors. This option may somewhat mitigate the issue of graph size explosion if the number of newly inserted vectors exceeds the number of deleted vectors, but may also lead to unpredictable degradation in the HNSW graph accuracy since edges are being modified. Emerging approaches~\cite{yu2025topology, xiao2024enhancing, xu2025place} show that lightweight edge-repair algorithms allow for complete deletions of nodes with minimal latency. \texttt{hnswlib} does not natively support these algorithms yet, and integration of edge repair into \SystemName{} is left as future work.
    
\section{Evaluation}
We evaluate \SystemName{} to answer several questions:
\begin{itemize}
    \item[\textbf{Q1.}] Does the partitioning co-locate nearest neighbors on the same partition, and does it remain effective as data is inserted and deleted? (Sections~\ref{sec:partitioning-quality} and~\ref{sec:partitining-updates})
    \item[\textbf{Q2.}] Does the routing layer with its three routing modes dispatch queries to the partitions containing their nearest neighbors? (Sections~\ref{sec:routing-quality} and~\ref{sec:routing-under-updates})
    \item[\textbf{Q3.}] Does \SystemName{} sustain high recall and throughput end-to-end on dynamic workloads that interleave insertions, deletions, and queries? (Section~\ref{sec:eval-perf})
    % \item[\textbf{Q4.}] Does the maintenance policy trigger rebuilds at the right points---often enough to prevent recall degradation, rarely enough to avoid wasted work? (Section~\ref{sec:eval-maintenance})
    % \item[\textbf{Q5.}] Does throughput scale with the number of workers while preserving recall? (Section~\ref{sec:eval-scalability})
\end{itemize}

\subsection{Experimental Setup}
\label{sec:eval-setup}
\textbf{Hardware.} All experiments run on CloudLab Utah \texttt{d6515} nodes \cite{Duplyakin+:ATC19}. Each node has a single-socket AMD EPYC 7452 (32 cores, 64 threads at 2.35\,GHz), 128\,GB of DRAM, and a 25\,Gbps Ethernet experiment link. We pin all worker and coordinator processes to the 32 logical CPUs to reduce scheduling noise. Distributed experiments use a cluster of ten workers and one coordinator; single-node experiments use one node.

\textbf{Datasets.} We use two billion-scale datasets: SIFT1B~\cite{jegou2011searching, simhadri2022results} (128-dimensional, $L_2$) and MSTuring1B~\cite{simhadri2022results} (100-dimensional, $L_2$).

\textbf{Workloads.} For static experiments we use 500M-vector subsets of each dataset together with the standard query sets. 

% \paragraph{Dynamic workloads.}
For dynamic experiments we generate runbooks using the streaming-track workload generator from the NeurIPS 2023 Big-ANN Challenge~\cite{simhadri2024results}, on both 100M and 500M subsets. Runbooks consist of five rounds, comprising $1{,}280$ batches in total for clustered and uniform variants and roughly $616$ for shifting variants. To produce \emph{clustered} runbooks, we partition the dataset into 64 clusters via $k$-means; each round visits every cluster---inserts
a sample from the cluster, runs the full query set, deletes a fraction of the cluster, and runs the query set again. \emph{Uniform} runbooks share the same outer structure but draw insertion and deletion samples uniformly from the full dataset rather than from clusters. \emph{Shifting} runbooks instead partition the 64 clusters into five disjoint cohorts, one per round; each round inserts the clusters of a new cohort and evicts those inserted in the previous round, so the active set's support migrates across rounds rather than remaining fixed. Clustered and shifting runbooks specifically stress the routing layer's ability to track distribution shift over time.

% For dynamic experiments we generate runbooks using the streaming-track workload generator from the NeurIPS 2023 Big-ANN Challenge~\cite{simhadri2024results}, on both 100M and 500M subsets. Runbooks consists of either five rounds of 256 batches, for 1{,}280 batches in total or . To produce \emph{clustered} runbooks, we partition the dataset into 64 clusters via $k$-means; each round visits every cluster---insert a sample from the cluster, run the full query set, delete a fraction of the cluster, and run the query set again. \emph{Uniform} runbooks share the same outer structure but draw insertion and deletion samples uniformly from the full dataset rather than from clusters. \emph{Shifting} runbooks \TODO{}. Clustered runbooks specifically stress the routing layer's ability to track distribution shift over time.

 \textbf{Baselines.} We compare against GP-ANN~\cite{gottensuren2025unleashing}, a distributed graph-based system over static data; Quake~\cite{mohoney2025quake}, a single-node dynamic partitioned index; and PostgreSQL with the \texttt{pgvector} and Citus extensions \cite{postgresql,pgvector,cubukcu2021citus} for distributed dynamic workloads. All baselines run on the same \texttt{d6515} hardware. Quake \cite{mohoney2025quake} is included only for recall comparison---a throughput comparison between a distributed system and a single-node baseline is not meaningful. For the routing graph produced by GP-ANN we use identical HNSW parameters. We implement a simple insertion script for GPANN using their existing distributed benchmarking script. Inserted vectors are routed to their closest partition and inserted into the HNSW graph for that partition.

\textbf{Metrics.} We report \emph{recall@$k$} (fraction of true top-$k$ neighbors retrieved) and \emph{throughput} (queries per second, QPS) as the primary end-to-end metrics. To isolate routing quality from in-partition search quality we additionally report \emph{theoretical recall@$k$} (Section~\ref{sec:sec:eval-metrics}), which assumes exhaustive search within each probed partition. We report \emph{partition activation rate}---the fraction of workers engaged per query---as a routing-efficiency metric, and \emph{partition imbalance}---the Gini coefficient of per-worker vector counts---as a load metric. 

\textbf{Routing Oracles.}
\label{sec:routing-oracles}
To separate the quality of the partitioning from the quality of the routing decision, we compare each routing mode against an \emph{oracle} that uses ground-truth neighbors to make the best possible partition selection under the same budget.

\textit{NProbe oracle.} For an nprobe value $N$ and query $q$, the oracle visits the $N$ partitions containing the largest fraction of $q$'s true nearest neighbors and reports the resulting recall. This is the highest recall any router can reach while probing exactly $N$ partitions, so it measures how well the partitioning co-locates neighbors. The gap between oracle recall and observed recall measures the amount of possible recall lost due to imperfect routing.

\textit{Routing Factor oracle.} % NOTE: this oracle is the least informative of the three (it is conditioned on the router's own partition count $n_q$); consider dropping it for space and keeping only the NProbe and Recall-Target oracles.
For a routing factor $K$ and query $q$, we first simulate the actual routing: the top-$K$ centroids retrieved from the routing HNSW map to $n_q$ unique partitions. The oracle then selects the $n_q$ partitions holding the most ground-truth-neighbors. This measures router quality at a matched partition search budget: a small gap between observed and oracle recall indicates that the HNSW router correctly identifies the right partitions given a routing budget.

\textit{Recall Target oracle.} For a recall target $\hat{r}$ and query $q$, the oracle ranks partitions by ground-truth-neighbor fraction and greedily accumulates them in that order until the cumulative recall reaches $\hat{r}$, reporting the mean number of partitions visited and the mean recall achieved. This is the minimum number of partitions any router needs to attain $\hat{r}$ on average, and the gap between the oracle and observed recall is a measure of how effectively the adaptive routing algorithm identifies partitions to reach the recall target without searching too much or too little.
\begin{figure*}
    \centering
    \includegraphics[width=1\textwidth]{figures/static_routing_quality.pdf}
    \caption{\textbf{Static routing quality across the three routing modes.} Each column is a dataset (SIFT-100M, MSTuring-100M, SIFT-500M, MSTuring-500M); each row is a routing mode. \emph{Top (NProbe):} theoretical recall@10 versus the number of partitions probed for \SystemName{}, its NProbe oracle, and GP-ANN (GP-ANN omitted on the 500M datasets, where its partition build timed out). \emph{Middle (Routing Factor):} theoretical recall@10 and partition activation versus routing factor $K$ for \SystemName{} and the Routing Factor oracle. \emph{Bottom (Recall Target):} theoretical recall@10 and partition activation versus the requested recall target for \SystemName{} and the Recall Target oracle; the dotted diagonal marks the target. Recall curves rise steeply and saturate well before all ten partitions are probed, indicating that each query's true neighbors are concentrated in a few partitions. \SystemName{} tracks the oracle closely on SIFT and matches GP-ANN despite partitioning only $10^4$ centroids rather than the full dataset; the larger \SystemName{}--oracle gap on MSTuring reflects that dataset's harder routing structure.}
    \label{fig:static-datasets}
\end{figure*}

\subsection{Static Index Quality}
We first evaluate \SystemName{}'s performance on static datasets to establish a performance baseline before introducing updates. In this section we show (1) whether the partitioning co-locates nearest neighbors on the same partition, and (2) whether the three routing modes can identify the correct partitions without over-searching. To isolate partitioning quality from local graph quality, all recall values in this section are reported as \emph{theoretical recall@10} (see~\ref{sec:eval-metrics}), which assumes exhaustive search within each probed partition and therefore measures an upper bound on what any in-partition search algorithm can achieve given the routing decision.

\subsubsection{Partitioning Quality}
\label{sec:partitioning-quality}
Figure \ref{fig:static-datasets} plots theoretical recall@10 as a function of the number of partitions probed for all four datasets, alongside the NProbe oracle that greedily selects the best possible set of partitions at each nprobe budget. The shape of the recall curve is itself diagnostic. A partitioning that scatters nearest neighbors uniformly across partitions would produce a linear recall curve, reaching 1.0 only when all ten partitions are probed. The concavity we observe — a steep initial rise that saturates well before nprobe = 10 — reveals that the true nearest neighbors of most queries are concentrated in a small number of partitions.

At nprobe = 1, dispatching a query to a single partition, \SystemName{} already recovers approximately 70\% of true top-10 neighbors on SIFT subsets and approximately 65\% on MSTuring subsets. By nprobe = 3 on SIFT and nprobe = 6 on MSTuring, theoretical recall exceeds 95\%. The 500M curves follow similar trends, indicating that both the centroid-based partitioning and the routing HNSW scale to $5\times$ more data without degradation. 

The gap between \SystemName{} and the NProbe oracle at each nprobe quantifies the gap left by imperfect routing — the recall loss attributable to the router directing queries to the wrong partitions even when the partitioning itself is good. This gap is small on SIFT (under 5\% at nprobe = 3) and wider on MSTuring (roughly 10\% at nprobe = 3), pointing to a harder routing problem on the latter dataset rather than a failure of the underlying partition. We examine the dataset trends further in Section~\ref{sec:partitining-updates}.

Figure~\ref{fig:static-datasets} also shows GP-ANN's NProbe recall curve on the 100M datasets (building GP-ANN partitions for the 500M datasets timed out on our hardware). \SystemName{} achieves theoretical recall comparable to GP-ANN despite a much coarser partitioning strategy: GP-ANN runs KaHIP directly on an approximate KNN graph over the full dataset, while \SystemName{} partitions a routing HNSW over only $10^4$ centroids. That the two produce similar partition quality is consistent with the finding in~\cite{gottensuren2025unleashing}, although their end-to-end performance deviates. 

% Figure \ref{fig:static-datasets} reports theoretical recall@10 for the three routing modes on SIFT-100M, SIFT-500M, MSTuring-100M, and MSTuring-500M, together with the per-mode oracles from Section \ref{sec:routing-oracles}; for the nprobe mode we also show the GP-ANN baseline ofr the 100M datasets (building GP-ANN partitions for the 500M datasets timed out on our hardware). For Routing Factor and Recall Target we also show the average partition activation, which is the fraction of partitions searched.

 % Across all four datasets, \SystemName{}'s recall curve rises steeply from nprobe = 1 and approaches saturation well before nprobe = 10. At nprobe = 1 --- equivalent to dispatching the query to a single worker --- SURGE recovers approximately $70\%$ of the true top-10 neighbors on SIFT subsets and approximately $65\%$ on MSTuring subsets. By nprobe = 3 for SIFT and nprobe = 6 for MSTuring, the recall exceeds 0.95 on the datasets. The trends of the 500M curves are essentially indistinguishable in shape from their 100M counterparts, indicating that both partitioning and routing do not lose performance as the index scales up by 5x.

% \SystemName{} achieves theoretical recall comparable to GP-ANN. This is consistent with the two systems' design choices. GP-ANN constructs an approximate KNN graph over the full base dataset and partitions that graph directly; SURGE partitions a small HNSW built over centroids representative of the base distribution. Both are able to co-locate neighbors effectively and have various performance tradeoffs \dk{TODO: rephrase. basically, we know that qps vs recall is better for GP-ANN}, but \SystemName{}'s routing layer is much easier to rebuild periodically in the case of distribution drift.


 % The concavity of the curves quantifies a critical property: for the typical query, the true nearest neighbors lie inside a subset of the ten partitions rather than being spread across all of them. A partitioning that did not co-locate neighbors --- for example, one that assigned vectors to workers by random hashing --- would produce a linear increase in theoretical recall, reaching 1.0 only at nprobe = 10. Any routing scheme over such a partitioning would have to probe every worker to deliver high recall, eliminating the throughput advantages of distribution.

 \subsubsection{Routing Quality}
 \label{sec:routing-quality}
 Given a fixed partitioning scheme, the routing layer must dispatch each query to the partitions most likely to contain its nearest neighbors, using only the routing HNSW and the centroid $\rightarrow$ partition mapping. Figure~\ref{fig:static-datasets} evaluates whether the three routing modes accomplish this, and how efficiently.

 Under NProbe mode, \SystemName{} tracks the oracle closely across all four datasets, confirming that the routing HNSW identifies the correct partitions at each probe budget. Under Routing Factor mode, the top-K centroids typically collapse to well below K distinct partitions, consistent with the partitioning co-locating proximate centroids. 

 Recall Target mode overshoots the specified targets on all datasets, which is expected: the oracle exhibits the same behavior, since neighbors are concentrated within partitions, probing the minimum number of partitions needed to reach the target recall consistently yields theoretical recalls above it. On SIFT, \SystemName{}'s recall and activation closely follow the oracle, indicating accurate partition selection. On MSTuring the gap to the oracle is larger, reflecting routing errors attributable to the dataset's more complex neighborhood structure, which causes the exponential distance weights to spread inconsistently across partitions. This suggests that more accurate recall-targeted routing could be achieved through dataset-specific calibration of the weighting scheme in Section~\ref{sec:dynamic-routing}.
 
 % Recall Target mode overshoots the recall targets for all datasets. However, the behavior is consistent with that of the oracle, which indicates the average recall achieved by ensuring that each query hits a recall $\geq$ the target. For the SIFT datasets, the recall target router's recall and activation closely follows the behavior of the oracle, indicating that the routing can accurately hit the recall target while correctly choosing the partitions to probe. For MSTuring, the router does not perform as accurately, indicating that it is choosing some of the wrong partitions to probe. These errors are a product of dataset complexity, which causes the exponential weights to spread inconsistently across partitions depending on the local neighborhood structure around each query in MSTuring. This indicates that the current routing is a strong approximating, and more accurate recall-targeted routing schemes can be produced by calibrating to the dataset.
 
 % Approaches that observe intermediate results before expanding the search, such as Quake's Adaptive Partition Scanning~\cite{quake}, avoid this failure mode but serialize partition probes; SURGE accepts this trade-off in exchange for fully parallel dispatch.

 
 % This section evaluates whether the routing layer dispatches queries to the partitions actually containing their nearest neighbors, without over-searching. 
 
 % Figure~\ref{fig:static-datasets} plots theoretical recall@10 and partition activation rate as a function of (a) routing factor $K$ and (b) recall target $\hat{r}$ on SIFT100M and MSTuring100M, comparing \SystemName{} against Quake's APS adaptive scanning~\cite{mohoney2025quake} in (b) (\TODO{}). For each recall target, we report the recall and activation achieved by the target recall oracle \ref{sec:routing-oracles}. The recall reported in the graph is the average across all queries for each recall target.

 \subsubsection{Search Performance}
 \begin{figure}
    \centering
    \includegraphics[width=1\linewidth]{figures/static_qps_combined.pdf}
    \caption{\textbf{End-to-end recall--throughput on static datasets.} For each dataset (one panel each, SIFT-100M, MSTuring-100M, SIFT-500M, MSTuring-500M), query throughput (QPS, log scale) and partition activation versus recall@10 for \SystemName{}'s three routing modes and GP-ANN, all on the same ten-worker cluster. Build and HNSW parameters are fixed; only the routing parameter (nprobe, routing factor, or recall target for \SystemName{}; nprobe for GP-ANN) is varied to trace each curve. Routing Factor mode attains slightly higher throughput than NProbe at equal recall because it dispatches after a single routing search; Recall Target mode is comparable to Routing Factor. \TODO{The SIFT-100M and SIFT-500M GP-ANN curves are being re-run (GP-ANN currently underperforms \SystemName{}, which is inconsistent with~\cite{gottensuren2025unleashing}); update this caption and the surrounding text once the expected GP-ANN advantage is reproduced.}}
    \label{fig:perf-static}
\end{figure}

% Figure~\ref{fig:static-search-perf} plots end-to-end recall@10 versus throughput for SURGE's three routing modes and GP-ANN on all four datasets, varying only the routing parameter to trace the Pareto frontier. 

Figure~\ref{fig:perf-static} plots end-to-end recall@10 versus query throughput (QPS) on all four datasets for \SystemName{}'s three routing modes and GP-ANN, all running on the same ten-worker cluster. Build and HNSW parameters are held constant; we vary only the routing parameter (nprobe, routing factor, or recall target for \SystemName{}; nprobe for GP-ANN) to trace the recall–throughput tradeoff for each system.

\dk{SIFT-100M - rerun b/c gp-ann is performing worse than us. Unexpected given the results in their paper. SIFT-500M rerun b/c at .7 recalll, routing factor is not following the trend nicely. Rerun both MSTurings and try to figure out why the recall target mode is doing so poorly at the lower recalls - was the weight equation there updated?}

Routing Factor mode achieves slightly higher throughput than NProbe mode at equal recall across all datasets. NProbe mode must re-query the routing HNSW with progressively larger K until it accumulates the required number of distinct partitions, whereas Routing Factor mode issues a single routing search and dispatches immediately to whatever partitions the top-K centroids map to. Recall Target mode achieves comparable performance to Routing Factor. For MSTuring, all modes converge to the same throughput at high recall because they require searching all or nearly all partitions.

% GP-ANN's recall–throughput curve lies to the right of SURGE on both datasets: at equal recall, GP-ANN delivers higher QPS on static workloads. GP-ANN demonstrates in~\cite{gottesbueren2025} that graph partitioning outperforms the centroid-based partitioning (Pyramid~\cite{deng2019pyramid}), on which SURGE's routing architecture is based.The 500M curves preserve the same shape and relative ordering as the 100M results, confirming that neither routing quality nor throughput degrades as the index scales fivefold.


At equal recall, GP-ANN delivers \TODO{measure: X}$\times$ higher QPS than \SystemName{} on SIFT-100M and \TODO{measure: X}$\times$ on MSTuring-100M. GP-ANN demonstrates in~\cite{gottensuren2025unleashing} that their approximate KNN graph partitioning outperforms centroid-based partitioning (Pyramid~\cite{deng2019pyramid}), on which \SystemName{}'s routing architecture is based. \SystemName{} accepts this throughput gap in exchange for a partitioning that can be maintained online under updates.

% GP-ANN directly partitions an approximate KNN graph over all dataset vectors, yielding finer partition boundaries that require fewer partitions to be probed per query. SURGE accepts this throughput gap in exchange for a partitioning that can be maintained online under updates, the central capability evaluated in Section 4.3. The 500M results preserve the same relative ordering as the 100M curves, confirming that routing quality and throughput scale without degradation to 5× more data.

 % Figure~\ref{fig:perf-static} plots recall@10 vs.\ QPS on SIFT100M and MSTuring100M for \SystemName{}, GP-ANN, and PostgreSQL+pgvector+Citus (\TODO{}). All systems run on the same ten-worker cluster, and \SystemName{} reports performance for every routing mode (\S\ref{sec:dynamic-routing}). Build and HNSW parameters are held constant for every search, we vary only nprobe for GP-ANN or the routing parameter (nprobe, routing factor, recall target) for \SystemName{}, which allows us to isolate the performance of routing and the partitions rather than local or routing graph build parameters.

\subsection{Dynamic Index Quality}
\label{sec:eval-routing}



% \dk{Expected: \textbf{(Routing factor)} Recall increases with $K$ and then plateaus as recall reaches 100\%; activation rises sub-linearly because the partitioning co-locates centroids, so increasing $K$ often reveals more centroids already on the same partitions. \textbf{(Recall target)} \SystemName{}'s per-query partition selection should achieve the requested target with activation close to the "minimal oracle". Based on current experiments, I expect \SystemName{}'s recall-target mode to outperform Quake. Discussion will emphasize that \SystemName{} achieves the target recalls and makes the routing decision \emph{before} dispatch.}

\subsubsection{Partitioning Under Updates}
\label{sec:partitining-updates}
\begin{figure*}
    \centering
    \includegraphics[width=.7\textwidth]{figures/theoretical_runbook_grid_recalltarget.pdf}
    \caption{\textbf{Recall-target routing under updates.} Theoretical recall@10 (top sub-row) and partition activation (bottom sub-row) across runbook steps at a recall target of 0.9, for SIFT-100M (top row) and MSTuring-100M (bottom row) under uniform, clustered, and shifting update patterns (columns). Curves: \emph{fresh routing} (routing layer rebuilt each step), \emph{stale routing} (routing layer fixed from the first insert step), and the corresponding \emph{fresh}/\emph{stale} oracles. Shaded bands mark insert (green) and delete (pink) phases. On uniform workloads all curves coincide---uniform updates do not shift centroids. On clustered and shifting workloads the stale oracle stays high (the partition layout still co-locates neighbors) while stale routing falls well below it, isolating the degradation to the router rather than the partitions. The fresh oracle reaches the target with fewer activated partitions than the stale oracle, showing that maintenance produces tighter partitions rather than merely restoring the original ones.}
    \label{fig:theoretical-runbooks-recalltarget}
\end{figure*}

\begin{figure*}
    \centering
    \includegraphics[width=.7\textwidth]{figures/theoretical_runbook_grid_nprobe.pdf}
    \caption{\textbf{Fixed-nprobe routing under updates.} Theoretical recall@10 (top sub-row) and partition imbalance (bottom sub-row, identical across routing modes) across runbook steps for SIFT-100M at nprobe~$=3$ (top row) and MSTuring-100M at nprobe~$=5$ (bottom row), under uniform, clustered, and shifting updates (columns). Curves: \emph{fresh routing}, \emph{stale routing}, their \emph{fresh}/\emph{stale} oracles, and \emph{GP-ANN}, whose routing layer is likewise fixed from the first insert step. Shaded bands mark insert (green) and delete (pink) phases. Stale routing and GP-ANN degrade together on clustered and shifting workloads---confirming that routing degradation is a structural consequence of any fixed centroid representation, not an artifact of one partitioning scheme---while fresh routing tracks its oracle. Partition imbalance spikes at the onset of the first insert phase on non-uniform workloads and then declines as later rounds fill sparse partitions; \SystemName{}'s imbalance is somewhat higher than GP-ANN's because maintenance trades balance to co-locate neighbors.}
    \label{fig:theoretical-runbooks-nprobe}
\end{figure*}

% The oracle curves in Figures~\ref{fig:theoretical-runbooks-recalltarget} and \ref{fig:theoretical-runbooks-nprobe} isolate partition quality from routing quality: at each runbook step, the oracle selects the optimal set of partitions given ground-truth neighbor locations and reports the resulting recall. If the partitioning itself degrades under updates, the oracle recall will fall regardless of how the router behaves. Across all three runbook types and both datasets, both the fresh and stale oracles remain consistently high --- on SIFT-100M the stale NProbe oracle stays above 90\% throughout clustered and shifting workloads, and on MSTuring-100M above 95\%, despite neither being rebuilt. This indicates that the original partitioning continues to co-locate nearest neighbors as updates arrive and that degradation in the observed recall is attributable to the routing layer, not to the partition layout.

The oracle curves in Figures~\ref{fig:theoretical-runbooks-recalltarget} and \ref{fig:theoretical-runbooks-nprobe} isolate partition quality from routing quality: at each runbook step, the oracle selects the optimal set of partitions given ground-truth neighbor locations and reports the resulting recall. Both the fresh and stale oracles remain high across all runbook types and datasets — on SIFT-100M the stale NProbe oracle stays above 90\% on clustered and shifting workloads, and on MSTuring-100M above 95\%. Since oracle recall is the best any router can achieve given the current partition layout, the stability of the stale oracle indicates that the partitioning continues to co-locate nearest neighbors as updates arrive and that the recall degradation in \SystemName{}'s observed recall is attributable to the router, not to the partition layout.

% The gap between the two oracles is informative. In recall-target mode (Figure~\ref{fig:theoretical-runbooks-recalltarget}), the oracle built on a freshly maintained routing layer consistently activates a smaller fraction of partitions than the stale oracle to reach the same recall target, across all three runbook types and both datasets. In NProbe mode (Figure~\ref{fig:theoretical-runbooks-nprobe}), the fresh oracle achieves slightly higher theoretical recall at the same probe budget, \TODO{with the advantage widening on clustered and shifting workloads (verify on the msturing shifting workload) }. Both observations reflect the same underlying effect: maintenance produces tighter partitions that concentrate each query's nearest neighbors more compactly, reducing the number of partitions that must be visited to achieve a given recall. Even so, the stale oracle remains slightly above \SystemName{}'s observed recall lines on harder workloads and datasets (MSTuring-100M, Clustered and Shifting runbooks), indicating that the router does not produce ideal search choices, but rather a good approximation.

The fresh oracle consistently outperforms the stale oracle, indicating that maintenance improves---rather than merely recovers---partition quality. In recall-target mode (Figure~\ref{fig:theoretical-runbooks-recalltarget}), the fresh oracle activates a smaller fraction of partitions to reach the same recall target on all runbook types and both datasets. In NProbe mode (Figure~\ref{fig:theoretical-runbooks-nprobe}), this advantage appears as a consistent recall gain at the fixed probe budget, most pronounced on SIFT and on the clustered and shifting workloads. Both observations reflect the same property: fresh partitions concentrate each query's nearest neighbors more compactly, reducing the search budget required to cover them.

% The benefit of maintaining the routing layer varies significantly with workload type. On uniform runbooks, all four curves---fresh routing, stale routing, oracle fresh, oracle stale---are nearly indistinguishable throughout the full run. Uniform updates draw insertions and deletions from the entire dataset, leaving centroid positions largely unchanged and rendering maintenance unnecessary. Clustered and shifting runbooks produce a different outcome. On clustered workloads, the stale oracle remains high while stale routing falls below it, confirming that the partitions co-locate neighbors but the router can no longer identify which ones to probe. Shifting runbooks are the most demanding: the active dataset's distribution completely inserts and deletes entire groups of clusters in each round rather than cycling through a fixed set of clusters, so the router must continuously track an evolving distribution. The achieved recall for stale routing is therefore lower on shifting workloads than on clustered ones.

% The two datasets exhibit meaningfully different intrinsic routing difficulty, as already observed in Section\TODO{}. On SIFT-100M, an NProbe of 3 is sufficient to reach high oracle recall; MSTuring-100M requires NProbe 5 to achieve comparable quality under the same 10-worker setup. This structural gap carries over to the dynamic setting: the recall loss under stale routing is larger on MSTuring than on SIFT across all three workload types, since misrouting in a harder space is more costly.

The benefit of maintenance varies with workload type and dataset. On uniform runbooks, all four curves are indistinguishable throughout. Uniform updates do not shift centroid positions, and the partition layout and routing layer remains representative without maintenance. On clustered and shifting workloads, the oracle gap grows with distribution drift, with shifting runbooks being the more demanding case. MSTuring-100M is the harder dataset — requiring NProbe 5 to reach oracle recall comparable to SIFT-100M at NProbe 3 (Section~\ref{sec:partitioning-quality}) — and this structural difficulty amplifies the sensitivity to partition quality on harder runbooks.

% The non-oracle curves corroborate the partition quality interpretation. Fresh routing, which rebuilds the routing layer at every step, closely tracks the fresh oracle across all workloads, remaining within the oracle's recall range throughout. Even stale routing's activation (Figure 3) stays bounded: it rises modestly above the fresh routing activation on harder workloads but does not approach the maximum, indicating that the router does not compensate for degraded centroid positions by indiscriminately probing more partitions.

% The non-oracle routing lines are consistent with these trends. Fresh routing closely tracks the fresh oracle across all workloads, and even stale routing's partition activation rises only modestly above the stale oracle baseline on harder workloads, indicating that the router does not compensate for stale centroids through over-searching.

% Figures~\ref{fig:theoretical-runbooks-recalltarget} and \ref{fig:theoretical-runbooks-nprobe} reports theoretical recall@10 across runbook steps for both routing modes (fixed nprobe $K$ and fixed recall target), comparing \SystemName{} with a fresh routing layer on each step, \SystemName{} with a stale routing layer (router built on the first insert step), and GP-ANN naive insertions without maintenance (also using the routing layer built on the first insert step). Both include the respective oracle results (\ref{sec:routing-oracles}). The bottom panel in Recall Target mode shows the fraction of partitions searched by the different setups and the oracles during runbook execution. NProbe mode shows the imbalance across partitions (how many elements each partition contains), this value is the same for both modes.

% \dk{Expected: Without maintenance: recall falls (queries hit the wrong partitions) while activation rises (the router senses uncertainty and probes more). With maintenance, recall stays flat and close to the clean-build oracle. I expect that imbalance will be higher with lower partition activation since \TODO{}. GP-ANN clean builds will have the lowest imbalance but pay a full reindex cost at every step.}

\subsubsection{Routing Under Updates}
\label{sec:routing-under-updates}
% Fresh routing closely tracks the fresh oracle across all runbook types and both datasets, confirming that an updated routing layer is sufficient to direct queries to the partitions the oracle would select. On uniform runbooks, fresh and stale routing are indistinguishable, consistent with the absence of distribution shift. On clustered and shifting workloads the gap between them opens immediately at the onset of the first insert phase, as stale centroids lose alignment with the data distribution.

% Fresh routing closely tracks the fresh oracle across all workloads, and even stale routing's partition activation rises only modestly above the stale oracle baseline on harder workloads, indicating that the router does not compensate for stale centroids through over-searching. 

The gap between the oracle and the actual routing lines measures the quality of the routing decision independently of partition layout. Fresh routing--which uses a rebuilt routing layer after each update step in the runbook--tracks the corresponding oracle closely on all three workload types and both datasets. 

On uniform runbooks, fresh and stale routing are indistinguishable, consistent with the absence of distribution shift. On clustered and shifting workloads, the gap between them is apparent from the first insert phase and widens as stale centroids lose alignment with the evolving distribution.

% Recall loss is more substantial on harder workloads. Using nprobe routing mode (Figure~\ref{fig:theoretical-runbooks-nprobe}), stale routing on SIFT-100M falls to approximately 85\% on clustered workloads and reaches as low as 70\% in the early steps of shifting workloads, while the corresponding oracle remains above 90\%. On MSTuring-100M the degradation is more severe: using recall target mode (Figure~\ref{fig:theoretical-runbooks-recalltarget}), stale routing is unable to meet the target and remains at approximately 70--75\% on clustered runbooks and as low as 60--65\% on the shifting runbooks, a gap of up to 35 percentage points below the oracle.


Recall loss is more substantial on harder workloads. On SIFT-100M at NProbe 3 (Figure~\ref{fig:theoretical-runbooks-nprobe}), stale routing falls to approximately 85\% on clustered workloads and as low as 70\% in the early steps of shifting workloads, while the oracle remains above 90\%. On MSTuring-100M the degradation is more severe: at a recall target of 0.9 (Figure~\ref{fig:theoretical-runbooks-recalltarget}), stale routing settles at approximately 70--75\% on clustered runbooks and drops as low as 60--65\% on shifting runbooks, approximately $35$ percentage points below the oracle.

% GP-ANN, which also uses a static routing layer initialized from the first insert step, degrades alongside stale routing on all workloads with distribution shift. On clustered SIFT-100M it falls to approximately 90--92\% theoretical recall; on shifting MSTuring-100M to approximately 80--85\%. This confirms that routing degradation under distribution shift is a consequence of using a fixed centroid representation as data evolves, not a limitation of any particular routing implementation.

GP-ANN, which also builds and fixes its routing layer from the first insert step, degrades alongside stale routing on all workloads with distribution shift — reaching approximately 90--92\% on clustered SIFT-100M and 80--85\% on shifting MSTuring-100M with nprobe 3 and 5, respectively. That a separate system built on a different partitioning strategy exhibits the same trend confirms that routing degradation under distribution shift is a structural consequence of any fixed centroid representation, not a limitation of a specific implementation.

% Despite the recall loss, stale routing does not compensate by over-activating partitions. In recall-target mode, stale routing's activation fraction rises only modestly above the stale oracle baseline on harder workloads, indicating that the router probes the wrong partitions rather than probing more of them. In recall target mode this is expected, the routing decision is made purely on the assumption that the centroids are representative of the distribution, and if that is no longer the case the routing cannot be accurate.

Despite the recall loss, stale routing does not broaden its search to compensate. In recall-target mode, its activation fraction rises only modestly above the stale oracle baseline, indicating that misrouting manifests as probing the \emph{wrong} partitions rather than probing more of them. This is expected: as described in Section~\ref{sec:dynamic-routing}, partition scores are computed purely from centroid distances to the query, so once centroids are stale the router can no longer tell which partitions actually hold the neighbors.

% Partition imbalance (Figure~\ref{fig:theoretical-runbooks-nprobe}) follows a characteristic pattern on non-uniform workloads: a sharp spike at the onset of the first insert phase, as clustered batches concentrate vectors unevenly, followed by a steady decline as subsequent rounds fill previously sparse partitions. The imbalance is consistently higher than GP-ANN or stale partitioning because the maintenance operations trade some partitions balance in order to co-locate nearest neighbors.

Partition imbalance (Figure~\ref{fig:theoretical-runbooks-nprobe}) spikes sharply at the onset of the first insert phase on non-uniform workloads, as clustered batches concentrate vectors unevenly across partitions, then declines steadily as subsequent rounds distribute insertions across previously sparse partitions. The imbalance is consistently higher for \SystemName{} than GP-ANN or stale partitioning because the maintenance operations trade some partition balance in order to co-locate nearest neighbors. On uniform runbooks, imbalance remains negligible throughout.

\subsubsection{Maintenance Thresholds}
\begin{figure}
    \centering
    \includegraphics[width=1\linewidth]{figures/threshold_sweep.pdf}
    \caption{\textbf{Effect of the maintenance threshold $\tau$.} Theoretical recall@10 across runbook steps for $\tau \in \{0.0, 0.4, 0.6, 0.8, 1.0\}$ on MSTuring-100M (top) and SIFT-100M (bottom) at nprobe~$=3$, under clustered (left) and shifting (right) updates. $\tau=0$ rebuilds after every update batch; $\tau=1$ never rebuilds (stale routing). Crosses mark steps at which a rebuild is triggered. $\tau=1$ degrades progressively as centroid drift accumulates, whereas $\tau=0.6$ stays within $\approx$1\% of the $\tau=0$ recall on both datasets and both workloads while triggering far fewer rebuilds, and tightening to $\tau=0.4$ adds negligible recall. The same $\tau=0.6$ generalizes across datasets and workload types without per-dataset tuning, so it is adopted as the default in Section~\ref{sec:eval-perf}.}
    \label{fig:maintenance-threshold}
\end{figure}

Figure~\ref{fig:maintenance-threshold} isolates the effect of the maintenance trigger threshold $\tau$ by plotting theoretical recall@10 across runbook steps for $\tau \in \{0, 0.4, 0.6, 0.8, 1.0\}$ on clustered and shifting runbooks, the two workloads that stress the routing layer and require maintenance to maintain recall. $\tau = 0$ triggers maintenance after every update batch while $\tau = 1$ uses the initial, stale routing layer for the entire runbook. The steps during which maintenance is triggered are marked in the plot.

As expected, $\tau = 1$ causes recall to degrade progressively on both clustered and shifting runbooks as centroid drift accumulates, while $\tau = 0$ maintains recall - these are the same exact trends observed in Figure~\ref{fig:theoretical-runbooks-nprobe}. The more informative result is how little $\tau$ needs to be tightened to recover most of that recall. $\tau = 0.6$ — triggering a rebuild only when 60\% of centroids change partition assignment — achieves recall within 1\% of the $\tau=0$ case across both datasets and both runbook types, while reducing the number of rebuilds from 640 to \TODO{measure: rebuild count at $\tau{=}0.6$, clustered} on clustered runbooks and from 310 to \TODO{measure: rebuild count at $\tau{=}0.6$, shifting} on shifting runbooks. Tightening further to $\tau=0.4$ yields negligible additional recall improvement, while significantly increasing the number of maintenance operations.

Importantly, the same threshold generalizes across datasets and workload types without per-dataset tuning: the recall curves for $\tau=0.6$ track the $\tau=0$ baseline consistently on both SIFT and MSTuring, under both clustered and shifting update patterns. This suggests that the fraction of reassigned centroids is a reliable proxy for routing degradation regardless of the underlying data distribution. Based on these results, we adopt $\tau=0.6$ as the default for all end-to-end experiments in Section~\ref{sec:eval-perf}.

\subsection{End-to-End Performance on Runbooks}
\label{sec:eval-perf}

% \dk{Expected: \SystemName{} should be competitive with GP-ANN at matched recall, as both use HNSW routing over a partitioned dataset and incur no maintenance overhead on a static workload. PostgreSQL+Citus is expected to be substantially slower due to general-purpose query infrastructure overhead. The takeaway is that \SystemName{} is competitive in the static setting and that supporting updates costs little in terms of performance on static workloads.}


\subsubsection{Performance on Dynamic Workloads.}
\begin{figure*}
    \centering
    % \includegraphics[width=\textwidth]{figures/perf_dynamic.pdf} % TODO: generate and uncomment
    \fbox{\parbox[c][3.4cm][c]{0.95\textwidth}{\centering\textcolor{red}{\textbf{Placeholder: end-to-end dynamic performance.}}\\ \footnotesize recall@10, query QPS, and insert throughput vs.\ runbook step on SIFT-100M and MSTuring-100M (clustered, uniform, shifting) for \SystemName{} with maintenance, \SystemName{} without maintenance, \SystemName{} clean-rebuild-per-step, and GP-ANN clean-rebuild-per-step.}}
    \caption{\textbf{End-to-end performance on dynamic runbooks.} recall@10, query throughput, and insert throughput across runbook steps on SIFT-100M and MSTuring-100M under clustered, uniform, and shifting updates, for four configurations: \SystemName{} with maintenance ($\tau=0.6$), \SystemName{} without maintenance, \SystemName{} with a clean rebuild every step, and GP-ANN with a clean rebuild every step. GP-ANN is tuned to 95\% recall on step~1 and \SystemName{} uses a 95\% recall target. Steps at which \SystemName{} triggers maintenance are marked. \TODO{generate figure; confirm the configuration list and per-panel axis ranges match the final runs.}}
    \label{fig:perf-dynamic}
\end{figure*}

Figure~\ref{fig:perf-dynamic} reports recall@10, QPS, and insert throughput across runbook steps on MSTuring-100M and SIFT-100M (clustered, uniform, and shifting) for four configurations: \SystemName{} with maintenance, \SystemName{} without maintenance, \SystemName{} with a clean rebuild at every step (oracle), and GP-ANN with a clean rebuild at every step. To enable comparison, GP-ANN uses an nprobe factor which achieves 95\% recall on step 1, while \SystemName{} uses a recall target of 95\% to route queries. \SystemName{} uses a maintenance threshold of 60\%, and evaluates the system for maintenance after each step that contains a batch of updates. Steps where a maintenance operation occurs are marked.

\dk{Expected: \SystemName{} with maintenance should sustain roughly flat recall and QPS across the runbook, similar to the clean-build configurations. \SystemName{} without maintenance will show recall degradation on clustered runbooks while maintaining high QPS. Uniform runbooks should show smaller gaps between with/without maintenance, since uniform inserts do not shift centroids. Insert throughput: highest for no maintenance case, very low for clean rebuilds case. With maintenance should be reasonable, mark steps where the rebuilds occur.}

\subsubsection{End-to-End Time Breakdown}
\begin{table}
    \centering
    \caption{\textbf{End-to-end time breakdown on a fixed runbook.} Total wall-clock runtime and the percentage of time spent on search, insertion, deletion, and maintenance for \SystemName{}, Quake, and PostgreSQL+pgvector+Citus. Quake is single-node and is included to compare maintenance cost in a non-randomly-partitioned system. \TODO{populate with measured values.}}
    \label{tab:time-breakdown}
    \begin{tabular}{lrrr}
        \toprule
        & \SystemName{} & Quake & PostgreSQL+Citus \\
        \midrule
        Total (s) & \TODO{} & \TODO{} & \TODO{} \\
        Search    & \TODO{} & \TODO{} & \TODO{} \\
        Insert    & \TODO{} & \TODO{} & \TODO{} \\
        Delete    & \TODO{} & \TODO{} & \TODO{} \\
        Maint.    & \TODO{} & \TODO{} & \TODO{} \\
        \bottomrule
    \end{tabular}
\end{table}

Table~\ref{tab:time-breakdown} reports total wall-clock runtime and the percentage of time spent on each of search, insertion, deletion, and maintenance for \SystemName{}, Quake, and PostgreSQL+pgvector+Citus on a fixed runbook. Quake is a single-node system: we include it in the breakdown to compare maintenance costs in non-randomly partitioned systems.

\dk{Expected: \SystemName{} should spend the majority of its time on search and updates and a smaller fraction on maintenance. Quake may show smaller time for maintenance, but it can be attributed to being single node. Postgres should have the largest run time.}

% \subsection{Maintenance?}
% "is maintenance being performed at the correct points?"

% \subsection{Scalability?}
% Does throughput scale with number of workers while maintaining recall?


\section{Conclusion}

We presented \SystemName{}, a distributed approximate nearest neighbor search system that extends the advantages of HNSW \cite{malkov2018efficient} to billion-scale dynamic workloads. Prior distributed graph-based systems~\cite{deng2019pyramid, gottensuren2025unleashing} achieve high search quality on static datasets but require a full offline reindex to absorb any change to the underlying data, while existing dynamic ANN systems~\cite{singh2021freshdiskann, mohoney2025quake, xu2023spfresh} address graph maintenance only on a single node and offer no mechanism for keeping a distributed partitioning or routing layer current as data shifts. \SystemName{} closes this gap with three components. First, a neighborhood-preserving partitioning constructed over the base layer of a routing HNSW, maintained online through incremental centroid updates and refreshed by a re-partitioning and maximum-weight matching that minimizes cross-partition data movement during maintenance operations. Second, an adaptive maintenance policy that monitors centroid drift and triggers coordinated rebuilds only when the fraction of reassigned centroids exceeds a configurable threshold $\tau$, maintaining the ability to serve queries from existing indexes throughout the rebuild. Third, three routing modes: with a static routing factor, enforced n-probe, and a recall-targeted mode that decides which partitions to probe before dispatch. Alongside the system, we introduce theoretical recall as a metric that isolates routing quality from local search quality, and extend the NeurIPS '23 Big-ANN streaming generator \cite{simhadri2024results} to produce large scale runbooks on both clustered and uniform datasets, suitable for stressing update pipelines. 

Our evaluation on SIFT and MSTuring datasets and runbooks across a ten-worker cluster shows that \SystemName{} sustains \TODO{measure: recall\%} recall@10 at \TODO{measure: QPS} QPS over the full course of dynamic runbooks, achieving \TODO{measure: X}$\times$ the throughput of a distributed PostgreSQL with pgvector, while incurring \TODO{measure: X}$\times$ lower maintenance cost than a per-step reindex. On static workloads \SystemName{} matches GP-ANN within \TODO{measure: X\%} throughput at equal recall, confirming that adding support for updates has minimal effect on static index performance. We further show that recall under clustered updates is preserved and even outperforms that of complete reindex operations, whereas the same configuration without maintenance degrades by up to \TODO{measure: X\%} as routing centroids fail to capture the data distribution. Integrating lightweight edge repair for in-place deletions~\cite{yu2025topology, xiao2024enhancing, xu2025place} and the resumable routing-graph search discussed in~\cite{lu2026depth} are natural extensions we leave to future work.


%%
%% The acknowledgments section is defined using the "acks" environment
%% (and NOT an unnumbered section). This ensures the proper
%% identification of the section in the article metadata, and the
%% consistent spelling of the heading.
\begin{acks}

\end{acks}

%%
%% The next two lines define the bibliography style to be used, and
%% the bibliography file.
\bibliographystyle{ACM-Reference-Format}
\bibliography{sample-base}


%%
%% If your work has an appendix, this is the place to put it.
\appendix

% \section{Research Methods}

% \section{Online Resources}

\end{document}
\endinput
%%
%% End of file `sample-sigconf.tex'.
